\section{At most countably many multiple lowest windows}
\label{sec:spectral-generic-simplicity}

\begin{theorem}[AC-S45: simplicity outside a countable set of windows]
\label{AC-S45}
Let \(a(\ell)\) be the lowest eigenvalue of the actual
Weil-window operator on \((0,\ell)\), \(\ell>0\).
The set
\[
 \mathcal M=\{\ell>0:\dim\ker(A_\ell-a(\ell))>1\}
                                                               \tag{GS.1}
\]
is at most countable.  Moreover, \(a(\ell)\) is strictly decreasing.
Hence the actual lowest eigenvalue is
simple on a dense set of full Lebesgue measure, unconditionally.
The same statement holds for the parameter
\(\lambda=e^{\ell/2}>1\).

The proof gives precise one-sided information.  At every
\(\ell>0\), choose a real normalized lowest eigenvector
\(v\) with \(\tau_+(v)\ne0\), which exists by AC-S44.
Then
\[
 a(\ell+h)\leq a(\ell)-\tfrac12\tau_+(v)^2h+o(h)
                                      \quad(h\downarrow0).
                                                               \tag{GS.2}
\]
If \(\ell\in\mathcal M\), then
\[
                   a(\ell-h)-a(\ell)=o(h).               \tag{GS.3}
\]
This theorem does not show that \(\mathcal M\) is empty,
and it does not determine the parity of a simple minimum.
\end{theorem}
\begin{proof}
\proofstep{1. One full-line operator for nearby lengths.}
Fix \(\ell>0\) and \(h_0\in(0,\ell/4)\).
In (BT.4), retain all prime atoms with displacement at most
\(\ell+h_0\), and truncate the polar kernel only outside
\(|x|\leq\ell+h_0\).  The result is one fixed self-adjoint
translation-invariant operator
\[
 H=G+\sigma*,\qquad G=\tfrac12\log(1+D^2),
                     \quad\|\sigma\|_{\rm TV}<\infty.    \tag{GS.4}
\]
Its part on every interval of length at most \(\ell+h_0\)
is the corresponding actual window operator: atoms outside
the difference interval do not contribute, and the polar
cutoff does not change the compressed kernel.
Thus these nearby forms are restrictions of one form \(q_H\).
Their form domains are nested under zero extension, and
the variational principle gives
\[
                  a(\ell-h)\geq a(\ell)\geq a(\ell+h).
                                                               \tag{GS.5}
\]
The same argument for any two lengths shows that \(a\) is
nonincreasing on \((0,\infty)\).

\proofstep{2. A nonzero trace gives a linear decrease to the right.}
Write \(\alpha=a(\ell)\), let \(E\) denote zero extension,
and take the real normalized \(v\) in the statement.
By the established Carleman splitting, \(Ev\in\operatorname{Dom}(H)\).
So does \(S_hEv\), and their sum is an admissible form vector
on \((0,\ell+h)\).
Translation invariance gives
\(q_H(Ev)=q_H(S_hEv)=\alpha\).
For the cross term put
\[
 c_h=q_H(Ev,S_hEv)-\alpha\langle Ev,S_hEv\rangle.
\]
Inside \((0,\ell)\), \(HEv=\alpha v\).
The only remaining overlap with \(S_hEv\) is the exterior
strip \((\ell,\ell+h)\).  Hence
\[
 c_h=-\int_0^h C_v^+(s)v(\ell-h+s)ds+o(h)
                 =-\tau_+(v)^2h+o(h).                  \tag{GS.6}
\]
The first equality uses the exterior singular action in
(BT.6); its finite-measure part is bounded and is paired
with a function tending uniformly to zero on that strip.
The second equality is exactly the dominated product limit
proved in (BT.7)--(BT.10), with \(u=w=v\).

Strong continuity of whole-line translation on \(L^2\)
gives
\[
                       \|Ev+S_hEv\|_2^2\longrightarrow4.
\]
Its Rayleigh quotient is therefore
\[
 \alpha+\frac{2c_h}{\|Ev+S_hEv\|_2^2}
                   =\alpha-\tfrac12\tau_+(v)^2h+o(h).
                                                               \tag{GS.7}
\]
The variational principle proves (GS.2), with the exact
factor \(1/2\) supplied by the limiting denominator four.
For any \(\ell_2>\ell_1\), first choose a sufficiently small
\(h<\ell_2-\ell_1\) in (GS.2), then use monotonicity.
This gives \(a(\ell_2)\leq a(\ell_1+h)<a(\ell_1)\),
proving strict decrease.

\proofstep{3. A multiple minimum has a flat left difference quotient.}
Suppose now that the lowest eigenspace is multiple.
By (TM.2), choose a real normalized
\(u\in E_\alpha\cap H^1_0(0,\ell)\), with
\(u'\in E_\alpha\).  AC-S36 makes \(u'\) bounded.
Consequently the \(H^1\) representative of \(u\) is
Lipschitz, with
\[
                    |u(\ell-t)|\leq M t,
                       \qquad M=\|u'\|_\infty.          \tag{GS.8}
\]
Let
\[
 V_h=1_{(\ell-h,\ell)}Eu,
 \qquad U_h=Eu-V_h.
\]
We check that the sharp strip has small form energy, rather
than assuming that sharp restriction preserves an operator
domain.  Put
\(g_h(t)=h^{-1}u(\ell-h+ht)\) on \((0,1)\).
Its value and derivative are uniformly bounded by a constant
depending on \(M\).  For any fixed \(0<s<1/2\), its zero
extension has uniformly bounded \(H^s(\mathbb R)\) norm.
Indeed, the interior fractional integral is bounded using
\(|g_h(t)-g_h(r)|\leq M|t-r|\), and the two exterior
integrals are bounded by
\[
 C_s\int_0^1 |g_h(t)|^2
                         \bigl(t^{-2s}+(1-t)^{-2s}\bigr)dt< C.
                                                               \tag{GS.9}
\]
The integrability is exactly where \(s<1/2\) is used.
Scaling the \(L^2\) norm and fractional seminorm gives
\[
 \|V_h\|_2^2=O(h^3),\qquad
          \|V_h\|_{H^s}^2=O_s(h^{3-2s}).                \tag{GS.10}
\]
Since \(\tfrac12\log(1+\xi^2)\leq C_s(1+\xi^2)^s\)
and \(\sigma*\) is bounded on \(L^2\), it follows that
\[
             q_H(V_h)-\alpha\|V_h\|_2^2
                                      =O_s(h^{3-2s})=o(h).
                                                               \tag{GS.11}
\]
This also places \(U_h\) in the form domain of the smaller
window; no operator-domain assertion for \(U_h\) is required.

Because \(Eu\in\operatorname{Dom}(H)\) and
\(HEu=\alpha u\) on the support of \(V_h\), form polarization
and the disjoint decomposition \(Eu=U_h+V_h\) give exactly
\[
 q_H(U_h)-\alpha\|U_h\|_2^2
               =q_H(V_h)-\alpha\|V_h\|_2^2.             \tag{GS.12}
\]
The denominator \(\|U_h\|_2^2\) tends to one.
Equations (GS.11)--(GS.12), applied as a Rayleigh test on
the smaller interval, give \(a(\ell-h)\leq\alpha+o(h)\).
The reverse inequality in (GS.5) proves (GS.3).

\proofstep{4. An elementary countability argument.}
At a multiple point \(\ell\), choose a rational number
\[
                         -\tfrac12\tau_+(v)^2<b<0.
\]
Let \(F_b(x)=a(x)-bx\).
Equations (GS.2)--(GS.3) imply, for every sufficiently small
\(h>0\),
\[
             F_b(\ell-h)<F_b(\ell),
             \qquad F_b(\ell+h)<F_b(\ell).               \tag{GS.13}
\]
Thus \(\ell\) is a strict local maximum of \(F_b\).
For each fixed rational \(b\), the strict local maxima of
any real function on an interval are at most countable:
around each such point choose a rational-endpoint interval
contained in its strict-maximality neighborhood and containing
the point.  That point is the unique maximizer on the interval,
so a fixed rational interval can serve at most one such point.
There are only countably many pairs of rational endpoints.
Taking the union over rational \(b\) proves that \(\mathcal M\)
is countable.  This argument requires no unproved
eigenvalue differentiability or shape-derivative formula.
\end{proof}

\paragraph{What remains of the simple-even assertion.}
The theorem proves actual simplicity outside a countable
exceptional set; this is stronger than merely restating
simplicity as the Sobolev exclusion in AC-S44.
It leaves two separate questions: whether any exceptional
window exists, and whether every simple minimum is even.
No conclusion about the number or absence of parity switches,
or the growing-window theta/prolate selection, is asserted here.
